\section{Two Protocols, One Name}
\label{sec:ch14-two-protocols-one-name}
\index{WebSocket!subprotocol negotiation|(}%

The composed config says \code{GRAPHQL\_WEBSOCKET\_SUBPROTOCOL\_AUTO} and
chapter~\ref{ch:composition} could not ask what that meant. Answering it needs
two facts, and the first one is a piece of GraphQL history that catches
everybody once.

\index{graphql-transport-ws@\code{graphql-transport-ws}!versus \code{graphql-ws}}%
There are two WebSocket protocols for GraphQL. The old one came from Apollo's
library, \code{subscriptions-transport-ws}, and has been unmaintained for
years. The new one came from the \code{graphql-ws} library and is what
everything current implements. Their subprotocol names, the strings that go in
a \code{Sec-WebSocket-Protocol} header, are:

\begin{minted}[fontsize=\footnotesize,breaklines]{text}
graphql-ws            the old one, from subscriptions-transport-ws
graphql-transport-ws  the new one, from the graphql-ws library
\end{minted}

The library called \code{graphql-ws} speaks the protocol called
\code{graphql-transport-ws}, and the protocol called \code{graphql-ws} comes
from a library called something else. HotChocolate's constants file says so
plainly, and it is the only place in this stack that does:

\begin{minted}[fontsize=\scriptsize,breaklines]{csharp}
/// The sub-protocol name for the GraphQL over WebSocket Protocol
/// https://github.com/enisdenjo/graphql-ws/blob/master/PROTOCOL.md
public const string GraphQL_Transport_WS = "graphql-transport-ws";

/// The sub-protocol name for the GraphQL over WebSocket Protocol from Apollo
/// https://github.com/apollographql/subscriptions-transport-ws/blob/master/PROTOCOL.md
public const string GraphQL_WS = "graphql-ws";
\end{minted}

Carry that, because everything else in this section depends on it.

\subsection*{Offer both, and the two ends disagree}

\index{WebSocket!which end chooses}%
A WebSocket client lists the subprotocols it can speak and the server picks
one. So: open a connection offering both, in the order the router offers them,
and see what comes back. Against the Reviews subgraph, and against the router:

\begin{minted}[fontsize=\footnotesize,breaklines]{text}
offered:  graphql-transport-ws, graphql-ws

HotChocolate 16.6.0 on 5105  chose  graphql-ws
Cosmo Router 0.337.1 on 3002 chose  graphql-transport-ws
\end{minted}

The same list, in the same order, and the opposite answer. One of these two
processes is picking the older of the two protocols, and it is the one in front
of your domain code.

\index{WebSocket!registration order decides}%
The reason is not a preference anybody expressed. It is a loop:

\begin{minted}[fontsize=\scriptsize,breaklines]{csharp}
foreach (var protocolHandler in _protocolHandlers)
{
    if (webSocketManager.WebSocketRequestedProtocols.Contains(protocolHandler.Name))
    {
        _webSocket = await webSocketManager.AcceptWebSocketAsync(protocolHandler.Name);
        return protocolHandler;
    }
}
\end{minted}

The iteration is over the server's handlers and the client's list is only
tested for membership, so the client's ordering carries no weight at all. What
decides is the order the handlers were registered in, and that is:

\begin{minted}[fontsize=\footnotesize,breaklines]{csharp}
.AddApolloProtocol()                // graphql-ws, the old one
.AddGraphQLOverWebSocketProtocol(); // graphql-transport-ws, the new one
\end{minted}

Apollo first. A client offering both gets the old protocol, and a client
offering both is what the Cosmo Router is by default.

\subsection*{On the wire}

\index{WebSocket!the handshake, captured}%
None of that is visible from either end: the router logs nothing about a
successful upgrade and HotChocolate's Kestrel logs sit at \code{Warning}. So
the capture below came from putting a recording proxy in the middle, which is
\code{scripts/subscription-wire.mjs} in the companion repository, with Reviews
moved to another port so the proxy can take 5105.

One recording holds both cases, because the harness probes the subgraph before
the router subscribes to it. The probe, offering both:

\begin{minted}[fontsize=\footnotesize,breaklines]{text}
GET /graphql HTTP/1.1
sec-websocket-protocol: graphql-transport-ws, graphql-ws
user-agent: node

HTTP/1.1 101 Switching Protocols
Sec-WebSocket-Protocol: graphql-ws
\end{minted}

Then the router, on the configuration this chapter ends with:

\begin{minted}[fontsize=\footnotesize,breaklines]{text}
GET /graphql HTTP/1.1
Host: host.docker.internal:5105
User-Agent: Go-http-client/1.1
Connection: Upgrade
Sec-Websocket-Protocol: graphql-transport-ws
Sec-Websocket-Version: 13
Upgrade: websocket

HTTP/1.1 101 Switching Protocols
Sec-WebSocket-Accept: 0+wyIeBPin4+OloVC9wc4/bHeys=
Sec-WebSocket-Protocol: graphql-transport-ws
\end{minted}

and the frames that follow it:

\begin{minted}[fontsize=\tiny,breaklines]{text}
router ->   {"type":"connection_init","payload":{}}
subgraph -> {"type":"connection_ack"}
router ->   {"id":"d9t72ovgpr8s73bm6qs0","type":"subscribe","payload":{"query":
             "subscription($a: ID!){onReviewAdded(productId: $a){id rating author
             {__typename id}}}","variables":{"a":"UHJvZHVjdDoAAACgAAAAQIAAAAAAAAAW"},
             "extensions":{"initialPayload":{}}}}
subgraph -> {"id":"d9t72ovgpr8s73bm6qs0","type":"next","payload":{"data":
             {"onReviewAdded":{"id":"UmV2aWV3OkrunwGI8rlynU+zZQ2jRtY=","rating":4,
             "author":{"__typename":"Customer","id":"Q3VzdG9tZXI6AAAAwAAAAECAAAAAAAAABA=="}}}}}
router ->   {"id":"d9t72ovgpr8s73bm6qs0","type":"complete"}
\end{minted}

\index{WebSocket!the legacy message names}%
Two things to take from the frames. The subgraph's \code{next} payload contains
the customer's key and no name, which is
section~\ref{sec:ch14-what-an-event-costs}'s cost visible as bytes. And the
message names are \code{subscribe} and \code{complete}, which belong to the new
protocol. Take the pin off, leave the setting at auto, and the same capture
comes back with \code{start} and \code{stop} instead, which are the old
protocol's names for the same two messages. That is the whole visible
difference, and it is visible only here.

\subsection*{Which is why Mosaic names it}

\index{websocketSubprotocol@\code{websocketSubprotocol}!pinned in the composer input}%
The reviews entry in \code{federation/mosaic.yaml} gains three lines:

\begin{minted}[fontsize=\footnotesize,breaklines]{yaml}
    subscription:
      url: http://localhost:5105/graphql
      protocol: ws
      websocketSubprotocol: graphql-transport-ws
\end{minted}

The first two are defaults written out. The third is the one that changes
anything, and after it the composed config says
\code{GRAPHQL\_WEBSOCKET\_SUBPROTOCOL\_TRANSPORT\_WS} and the handshake offers
one subprotocol rather than two.

Both arrangements deliver events, so nothing was broken before this and nothing
is fixed by it. What was wrong was that the choice was being made by
registration order in somebody else's container, could not be seen from either
log, appears in no schema, and landed on the protocol its own ecosystem has
moved off. A default nobody chose and nobody can observe is worth one line of
configuration even when it works.

\index{spelling!of \code{websocketSubprotocol}}%
And now the part that cost me twenty minutes. That key is camelCase, in a file
where every other key uses an underscore, and wgc's own documentation for the
file spells it with one too. The documented spelling composes perfectly and
does nothing at all, which is the subject of the next section.

\index{WebSocket!subprotocol negotiation|)}%
